% ============================================================================
% 第 8 章 总结与展望
% ============================================================================
\chapter{总结与展望}

% ----------------------------------------------------------------------------
\section{研究工作总结}

本文围绕“虚拟问诊系统设计与实现”这一课题，针对医学问诊实训中病例资源有限、标准化病人成本较高、教师批改压力大、学生课后练习机会不足等问题，设计并实现了一套基于 Web 与大语言模型的问诊训练系统。系统面向学生、教师和管理员三类角色，覆盖病例管理、AI 问诊训练、诊疗方案提交、训练评分、学习分析、班级教学分析、用户管理、日志审计和容器化部署等功能。

需求分析阶段，本文梳理了学生反复训练、教师病例建设与批改、管理员系统治理等需求，并提出性能、安全、可用性和可维护性要求。总体设计阶段，系统采用用户层、接入层、应用层、AI 服务层和数据层的分层架构；前端基于 Vue 3、TypeScript、Element Plus 与 ECharts 实现，后端基于 Spring Boot、Spring Security、MyBatis-Plus、Redis、MinIO 与 MySQL 实现，AI 能力通过阿里云百炼 DashScope 兼容接口接入 Qwen 系列模型。

核心实现方面，本文重点完成 AI Agent 层设计。系统把单提示词调用拆分为 PatientAgent、DiagnosisAgent 与 FeedbackAgent 三类角色，分别负责患者扮演、诊断评估和治疗反馈；通过 RAG 将医学知识点注入上下文，通过 Redis Memory 保持多轮问诊连续性，通过临床工具返回病例中的体格检查、辅助检查和影像资料，减少模型编造客观结果的风险；结构化 JSON 评分和规则兜底则用于生成训练报告。

系统实现与测试方面，本文完成学生端问诊训练、病例选择、训练报告和学习分析，教师端病例管理、批改工作台和教学分析，管理员端用户管理、系统监控、日志和存储管理等模块。测试结果表明，系统能够支持完整训练闭环，普通业务接口响应较稳定，AI 对话可通过流式输出改善等待体验，核心安全机制和容器化部署流程均可正常运行。

% ----------------------------------------------------------------------------
\section{创新点}

本文的主要创新点如下。

第一，构建面向医学问诊训练的多 Agent 角色分工机制。系统将问诊阶段的 PatientAgent 与提交后的 DiagnosisAgent、FeedbackAgent 分离，使模型在不同阶段遵守不同边界，降低角色混淆和提前泄题风险。

第二，设计临床工具自动注入机制。体格检查、化验和影像等客观结果不由模型自由生成，而是在调用模型前识别检查类请求，并优先从病例结构化字段和医学影像附件中返回，使系统能够区分“患者主观叙述”和“客观检查资料”。

第三，实现轻量级 RAG 与 Redis Memory 的组合。系统根据病例科室和用户问题进行关键词检索，将相关知识点注入提示词；同时使用 Redis 保存最近多轮对话，保证训练过程的上下文连续性。该方案部署简单、可解释性较强，适合中小规模教学系统。

第四，形成 AI 评分与教师批改结合的反馈闭环。系统输出总分及问诊、诊断、治疗、沟通等维度评分，并保存薄弱点、错误模式和改进建议；教师可在批改工作台复核和补充评价，在提高效率的同时保留教师主导作用。

第五，完成从前端交互、后端服务到容器化部署的工程实现。系统实现了用户认证、角色管理、病例管理、训练记录、对象存储、日志审计和 Docker Compose 部署，具备较完整的软件系统形态。

% ----------------------------------------------------------------------------
\section{工作不足}

受时间、数据来源和部署条件限制，本文仍存在以下不足。

首先，病例数据和医学知识库规模有限。当前系统主要使用模拟病例和结构化知识点进行验证，覆盖科室、疾病类型和难度层级还不够丰富。后续需要结合真实课程和教师资源，持续扩充经专家审核的病例与知识点。

其次，AI 评分仍有不确定性。虽然系统要求模型输出结构化 JSON，并实现了解析失败后的规则兜底，但模型评分与临床教师评分之间仍可能存在偏差。当前系统尚未基于大规模人工标注数据进行校准，因此 AI 评分更适合作为辅助评价，不能替代教师判断。

再次，临床工具自动注入仍以关键词和规则匹配为主。该方式简单可控，但对口语化、隐含式检查请求的识别能力有限。后续可结合意图识别模型或更细粒度的工具调用策略提高命中率。

最后，性能测试和安全测试仍处于毕业设计验证级别。普通接口、AI 对话和部署流程已经验证，但尚未经历真实大规模并发、长期运行、渗透测试和医疗数据合规审查。若进入实际教学环境，还需要加强压力测试、日志监控、备份恢复和隐私保护。

